% Dokumentum konfiguráció
\documentclass[fleqn,12pt]{article}
% --------------------------------------------------
% Dokumentum adatok
% --------------------------------------------------

% Intézmény és kar
\newcommand\INTEZMENY{Budapesti Műszaki és Gazdaságtudományi Egyetem}
\newcommand\KAR{Gépészmérnöki Kar}
\newcommand\LOGO{bme-logo.jpg} % logo fájl neve a figures mappában

% Dokumentum adatok
\newcommand\CIM{Festigkeitslehre Hausaufgabe}
\newcommand\DATUM{2025/26/1}
\newcommand\NEV{Abos Gergő}
\newcommand\NEPTUN{C46M9M}
\newcommand\EMAIL{gergo.abos@gmail.com}
\newcommand\TANTARGY{Festigkeitslehre}
\newcommand\TANTARGYKOD{BMEGEMMBXN2}

% --------------------------------------------------
% Preambulum és egyéni parancsok betöltése
% --------------------------------------------------

\input{../_settings/00_preambulum.tex}
\input{../_settings/01_commands.tex}
\input{../_settings/02_mechTikz.tex}

% ==================================================
% Dokumentum kezdete
% ==================================================
\begin{document}

% Címlap és tartalomjegyzék
% \input{contents/01_title.tex}
\includepdf{contents/Aufgabenblatt.pdf} % feladatlap beszúrása, ha szükséges
\section{Reaktionskrafte}
Erst ist die maßstäbliche Abbildung des Konstruktions zu sehen:
\begin{figure}[H]
  \centering
  \begin{tikzpicture}[thick]
    \def\mscale{3};
    \def\a{0.5 * \mscale};
    \def\b{2.0 * \mscale};
    \def\c{2.0 * \mscale};
    \def\h{0.433 * \mscale};

    \def\kNscale{0.5};
    \def\G{3 * \kNscale};

    \def\kNmscale{1.25};
    \def\p{0.5 * \kNmscale};

    \coordinate (A) at (-\c, -\h);
    \coordinate (B) at (\a / 2, -\h);
    \coordinate (C) at (0, 0);
    \coordinate (D) at (\a, 0);
    \coordinate (E) at (\a + \b, 0);

    \draw[vec={red}] (E) -- ++(0, -\G) node[left] {$G$};

    \fill[pattern=vertical lines, pattern color=red] (C) rectangle ($(E) + (0, \p)$);
    \draw[red] (C) rectangle ($(E) + (0, \p)$);
    \draw[force={red}] ($(E) + (0, \p)$) -- (E) node[midway, right] {$p$};

    \draw (A) -- (C);
    \draw (D) -- (C);
    \draw (B) -- (C);
    \draw (D) -- (B);
    \draw (D) -- (E);
    \draw [dashed, thin, color=gray] (A) -- (B);
    \supporthinge{A}{0}
    \supporthinge{B}{0}
    \hinge{C}
    \draw (D) ++(0.175, 0) arc (0:180:0.175);
    \hinge{D}
  \end{tikzpicture}
  \caption{Maßstäblich korrekte Abbildung des Konstruktions}
\end{figure}
Um die Reaktionen im Punkt A und B zu bestimmen, kann die Konstruktion als ein Starrkörper angenommen werden:
\begin{figure}[H]
  \centering
  \begin{tikzpicture}[thick]
    \def\mscale{3};
    \def\a{0.5 * \mscale};
    \def\b{2.0 * \mscale};
    \def\c{2.0 * \mscale};
    \def\h{0.433 * \mscale};

    \def\kNscale{0.5};
    \def\G{3 * \kNscale};

    \def\kNmscale{1.25};
    \def\p{0.5 * \kNmscale};

    \draw[force={red}] (E) -- ++(0, -\G) node[left] {$G$};

    \fill[pattern=vertical lines, pattern color=red] (C) rectangle ($(E) + (0, \p)$);
    \draw[red] (C) rectangle ($(E) + (0, \p)$);
    \draw[force={red}] ($(E) + (0, \p)$) -- (E) node[midway, right] {$p$};

    \draw (A) -- (C);
    \draw (D) -- (C);
    \draw (B) -- (C);
    \draw (D) -- (B);
    \draw (D) -- (E);

    \draw[force={red}, <-] (A) -- ++(-1, 0) node[midway, above] {$A_X$};
    \draw[force={red}, <-] (A) -- ++(0, -1) node[left] {$A_Y$};
    \draw[force={red}, <-] (B) -- ++(-1, 0) node[midway, above] {$B_X$};
    \draw[force={red}, <-] (B) -- ++(0, -1) node[left] {$B_Y$};

  \end{tikzpicture}
  \caption{Abbildung der Reaktionskräfte}
\end{figure}

Man bekommt 3 Gleichgewichtsgleichungen, und 1 Gleichung weil Stab 5 nur Kräfte nur im Stabrichtung auswirken kann.

\begin{align}
  &\sum F_X = A_X + B_X = 0 \\
  &\sum F_Y = A_Y + B_Y - G - p \cdot (a + b) = 0 \\
  &\sum M_A = B_Y \cdot (c + \frac{a}{2}) - G \cdot (a + b + c) - p \cdot (a + b) \cdot (c + \frac{a + b}{2}) = 0\\
  &\frac{A_X}{A_Y} = \frac{c}{\frac{a}{2} \cdot tg(60^\circ)}
\end{align}

Daraus bekommt man die folgende Werte für die Reaktionskräfte:
\begin{align*}
  \boxed{
    \begin{aligned}
      A_X &= -16{,}422 \, \text{kN} \\
      A_Y &= -3{,}556 \, \text{kN}
  \end{aligned}}
  &\qquad
  \boxed{
    \begin{aligned}
      B_X &= 16{,}422 \, \text{kN} \\
      B_Y &= 7{,}806 \, \text{kN}
  \end{aligned}}
\end{align*}

\section{Belastung der einzelnen Stäbe}
\subsection{Stab 5}
Da Stab 5 belastungen nur an seiner Endpunkte hat, sind die 2 Kräfte in diesen Punkte im Stabrichtung gerichet, und sind entgegengesetzt gleich groß.
\begin{figure}[H]
  \centering
  \begin{tikzpicture}[thick]
    \draw (C) -- (A);
    \draw[force] (C) -- ($ (C)!-1cm!(A) $) node[above] {$S_5$};
    \draw[force] (A) -- ($ (A)!-1cm!(C) $) node[above] {$S_5$};
    \draw[force, color=gray!50] (A) -- ++(1, 0) node[right] {$A_X$};
    \draw[force, color=gray!50] (A) -- ++(0, 1) node[right] {$A_Y$};
  \end{tikzpicture}
  \caption{Belastung von Stab 5}
\end{figure}
\begin{align*}
  |S_5| = |\mathbf{A}| = \sqrt{A_X^2 + A_Y^2} = 16,802 \, \text{kN}
\end{align*}
Da $A_X$ und $A_Y$ auch negativ sind:
\begin{align*}
  \boxed{
    S_5 = 16,802 \, \text{kN}
  }
\end{align*}

\subsection{Stab 4}
Nach dem selben Logik als Stab 5:
\begin{align*}
  &S_4 = G \\
  &\boxed{
    S_4 = 3 \, \text{kN}
  }
\end{align*}

\subsection{Stab 2 und 3}
Im Gelenk B sind 3 Kräfte zu finden. Die Summe der 3 Kräften muss 0 sein, für die Konstruktion im Gleichgewicht zu sein.
\begin{figure}[H]
  \centering
  \begin{tikzpicture}[thick]
    \draw[force, color=gray!50] (B) -- ++(1, 0) node[right] {$B_X$};
    \draw[force, color=gray!50] (B) -- ++(0, 1) node[above] {$B_Y$};
    \draw[dashed, color=gray!50] ($(B) + (1, 0)$) -- ++(0, 1);
    \draw[dashed, color=gray!50] ($(B) + (0, 1)$) -- ++(1, 0);
    \draw[force] (B) -- ++(1, 1) node[above right] {$\mathbf{B}$};

    \draw[force] (B) -- ++(-60:1.4) node[right] {$\mathbf{S_3}$};
    \draw[force] (B) -- ++(-120:1.4) node[left] {$\mathbf{S_2}$};
    \hinge {B};
  \end{tikzpicture}
  \caption{Belastung von Gelenk B}
\end{figure}

\setcounter{equation}{0}
\begin{align}
  \sum F_X = B_X + S_3 \cdot cos(-60^\circ) + S_2 \cdot cos(-120^\circ) = 0\\
  \sum F_Y = B_Y + S_3 \cdot sin(-60^\circ) + S_2 \cdot sin(-120^\circ) = 0
\end{align}
\begin{align*}
  \boxed{
    \begin{aligned}
      S_2 &= -20{,}928 \, \text{kN}
  \end{aligned}}
  &\qquad
  \boxed{
    \begin{aligned}
      S_3 &= 11{,}916 \, \text{kN}
  \end{aligned}}
\end{align*}

\section{Schnittgrößen von Stab 1}
\begin{figure}[H]
  \centering
  \begin{tikzpicture}[thick]
    \def\mscale{3};
    \def\a{0.5 * \mscale};
    \def\b{2.0 * \mscale};
    \def\c{2.0 * \mscale};
    \def\h{0.433 * \mscale};

    \def\kNscale{0.5};
    \def\G{3 * \kNscale};

    \def\kNmscale{1.25};
    \def\p{0.5 * \kNmscale};

    \draw[vec={red}] (E) -- ++(0, -1) node[left] {$G$};

    \fill[pattern=vertical lines, pattern color=red] (C) rectangle ($(E) + (0, \p)$);
    \draw[red] (C) rectangle ($(E) + (0, \p)$);
    \draw[force={red}] ($(E) + (0, \p)$) -- (E) node[midway, right] {$p$};

    \draw (D) -- (C);
    \draw (D) -- (E);

    \draw[force] (C) -- ($ (C)!1cm!(A) $) node[below] {$S_5$};

    \draw[force] (C) -- ($ (C)!1cm!(B) $) node[left] {$S_3$};
    \draw[force] (D) -- ($ (D)!1cm!(B) $) node[right] {$S_2$};

  \end{tikzpicture}
  \caption{Freikörperbild von Stab 1}
\end{figure}

Alle Kräfte werden zerlegt in x und y Komponente:
\begin{align*}
  \mathbf{S_5} &=
  \begin{pmatrix} S_{5X} \\ S_{5Y}
  \end{pmatrix}
  =
  \begin{pmatrix} A_X \\ A_Y
  \end{pmatrix}
  =
  \begin{pmatrix} -16{,}422 \\ -3{,}556
  \end{pmatrix} \text{kN} \\[2ex]
  \mathbf{S_3} &=
  \begin{pmatrix} S_3 \cdot \cos(-60^\circ) \\ S_3 \cdot \sin(-60^\circ)
  \end{pmatrix}
  =
  \begin{pmatrix} 11{,}916 \cdot 0{,}5 \\ 11{,}916 \cdot \left(-\frac{\sqrt{3}}{2}\right)
  \end{pmatrix}
  =
  \begin{pmatrix} 5{,}958 \\ -10{,}320
  \end{pmatrix} \text{kN} \\[2ex]
  \mathbf{S_2} &=
  \begin{pmatrix} S_2 \cdot \cos(-120^\circ) \\ S_2 \cdot \sin(-120^\circ)
  \end{pmatrix}
  =
  \begin{pmatrix} -20{,}928 \cdot (-0{,}5) \\ -20{,}928 \cdot \left(-\frac{\sqrt{3}}{2}\right)
  \end{pmatrix}
  =
  \begin{pmatrix} 10{,}464 \\ 18{,}124
  \end{pmatrix} \text{kN}
\end{align*}

Wegen eine Belastung im Punkt D werden die Schnittgrößenfunktionen aus 2 Teile bestehen:
\begin{enumerate}
  \item $0 < x < a$
  \item $a < x < a + b$
\end{enumerate}
\subsection{Teil 1. (0 < x < a)}
\begin{align*}
  N(x) &= S_{5X} + S_{3X} = -10,464 \, \text{[kN]} \\
  V(x) &= S_{5Y} + S_{3Y} - p \cdot x = -13,876 - 0,5x \, \text{[kN]} \\
  M_b(x) &= (S_{5Y} + S_{3Y}) \cdot x + p \cdot \frac{x^2}{2} = 13{,}876x + 0{,}25x^2 \, \text{[kNm]}
\end{align*}
\subsection{Teil 2. (a < x < a + b)}
\begin{align*}
  N(x) &= 0 \, \text{[kN]} \\
  V(x) &= G + p \cdot (a + b - x) = 4,25 - 0,5x \, \text{[kN]}\\
  M_b(x) &= G \cdot (a + b - x) + p \cdot \frac{(a + b - x)^2}{2} = \frac{x^2}{4} - 4,25x + 9,0625 \, \text{[kNm]}
\end{align*}

\subsection{Schnittgrößendiagramme}
\begin{figure}[H]
  \centering
  \begin{tikzpicture}[thick]
    \def\mscale{3};
    \def\a{0.5 * \mscale};
    \def\b{2.0 * \mscale};
    \def\c{2.0 * \mscale};
    \def\h{0.433 * \mscale};

    \def\kNscale{0.5};
    \def\G{3 * \kNscale};

    \def\kNmscale{1.25};
    \def\p{0.5 * \kNmscale};

    \draw[vec={red}] (E) -- ++(0, -1) node[left] {$G$};

    \fill[pattern=vertical lines, pattern color=red] (C) rectangle ($(E) + (0, \p)$);
    \draw[red] (C) rectangle ($(E) + (0, \p)$);
    \draw[force={red}] ($(E) + (0, \p)$) -- (E) node[midway, right] {$p$};

    \draw (D) -- (C);
    \draw (D) -- (E);

    \draw[force] (C) -- ($ (C)!1cm!(A) $) node[below] {$S_5$};

    \draw[force] (C) -- ($ (C)!1cm!(B) $) node[left] {$S_3$};
    \draw[force] (D) -- ($ (D)!1cm!(B) $) node[right] {$S_2$};

    \foreach \X in {0, \a, \a+\b} {
      \draw[dashed, gray] (\X, 0) -- (\X, -19);
    }

    %%%%%%%%%%%%%%%%%%%%%%%%%%%%%%%%
    % N(x)
    %%%%%%%%%%%%%%%%%%%%%%%%%%%%%%%%
    \begin{scope}[yshift=-3cm]
      \draw[->] (0,0) -- (\a + \b + 0.5, 0) node[below] {$x$};
      \draw[->] (0,-3) -- (0, 1) node[left] {$N(x)$} node[right] {$[kN]$};

      \def\yscale{0.25}
      % Step shape for N
      \draw[red, line width=1.5pt]
      (0, -10.464 * \yscale) --
      (\a, -10.464 * \yscale) --
      (\a, 0) -- (\a + \b, 0);

      % % Shaded area
      \fill[pattern=vertical lines, pattern color=red]
      (0,0) -- plot[domain=0:\a] (\x, { -10.464 * \yscale })
      -- (\a,0);

      % Labels
      \node[right, red] at (\a, -10.464 * \yscale ) {$-10.464$};
      \node[below left, red] at (\a+ \b, 0) {$0$};
    \end{scope}

    %%%%%%%%%%%%%%%%%%%%%%%%%%%%%%%%
    % V(x)
    %%%%%%%%%%%%%%%%%%%%%%%%%%%%%%%%
    \begin{scope}[yshift=-9cm]
      \draw[->] (0,0) -- (\a + \b + 0.5, 0) node[below] {$x$};
      \draw[->] (0,-3) -- (0, 1) node[left] {$V(x)$} node[right] {$[kN]$};

      \def\yscale{0.2}

      \draw[red, line width=1.5pt]
      (0,-13.876 * \yscale) -- plot[domain=0:\a] (\x, { (-13.876 - 0.5 * \x) * \yscale })
      -- (\a, -14.376 * \yscale);
      \fill[pattern=vertical lines, pattern color=red]
      (0,0) -- plot[domain=0:\a] (\x, { (-13.876 - 0.5 * \x) * \yscale })
      -- (\a,0);

      \draw[red, line width=1.5pt]
      (\a,-14.376 * \yscale) -- plot[domain=\a:\a+\b] (\x, { (4.25 - 0.5 * \x / \mscale) * \yscale })
      -- (\a+\b, 0 * \yscale);
      \fill[pattern=vertical lines, pattern color=red]
      (\a,0 * \yscale) -- plot[domain=\a:\a+\b] (\x, { (4.25 - 0.5 * \x / \mscale) * \yscale })
      -- (\a+\b, 0 * \yscale);

      % Labels
      \node[right, red] at (\a, -14.376 * \yscale ) {$-14.376$};
      \node[left, red] at (0, -13.876 * \yscale) {$-13.876$};
      \node[left, red] at (\a, 4 * \yscale) {$4$};
      \node[right, red] at (\a + \b, 3 * \yscale) {$3$};
    \end{scope}

    %%%%%%%%%%%%%%%%%%%%%%%%%%%%%%%%
    % M_b(x)
    %%%%%%%%%%%%%%%%%%%%%%%%%%%%%%%%
    \begin{scope}[yshift=-19cm]
      \draw[->] (0,0) -- (\a + \b + 0.5, 0) node[below] {$x$};
      \draw[->] (0,-1) -- (0, 5) node[left] {$M_b(x)$} node[right] {$[kNm]$};

      \def\yscale{0.6}

      \draw[red, line width=1.5pt]
      plot[domain=0:\a] (\x, { (13.876 * \x / \mscale - 0.25 * (\x / \mscale)^2) * \yscale });
      \fill[pattern=vertical lines, pattern color=red]
      (0,0) -- plot[domain=0:\a] (\x, { (13.876 * \x / \mscale - 0.25 * (\x / \mscale)^2) * \yscale })
      -- (\a,0);

      \draw[red, line width=1.5pt]
      plot[domain=\a:\a+\b] (\x, { ((\x / \mscale)^2/4 - 4.25 * \x / \mscale + 9.0625) * \yscale });
      \fill[pattern=vertical lines, pattern color=red]
      (\a,0 * \yscale) -- plot[domain=\a:\a+\b] (\x, { ((\x / \mscale)^2/4 - 4.25 * \x / \mscale + 9.0625) * \yscale })
      -- (\a+\b, 0 * \yscale);

      % Tangenten
      \draw[blue, dashed]
      (0,0 * \yscale) -- (\a/2, 3.43775 * \yscale) -- (\a, 7 * \yscale);
      \draw[blue, dashed]
      (\a,7 * \yscale) -- (\a+\b/2, 3*\yscale)
      -- (\a+\b, 0 * \yscale);
      \fill[blue] (\a/2, 3.43775 * \yscale) circle (0.05cm);
      \fill[blue] (\a+\b/2, 3 * \yscale) circle (0.05cm);
      \draw[black, dashed] (\a/2, 3.43775 * \yscale) -- (0, 3.43775 * \yscale) node[left]{$3,43775$};
      \draw[black, dashed] (\a+\b/2, 3 * \yscale) -- ++(1, 0) node[right]{$3$};
      % Labels
      \node[left, red] at (\a, 7 * \yscale ) {$7$};
    \end{scope}
  \end{tikzpicture}
\end{figure}

\section{Dimensionierung für reine Biegung}
Von den Schnittgrößen kann abgelesen werden:
\begin{align*}
  &M_{b,max}^{(1)} = 7kNm
\end{align*}
Nach dem Navier Formel:
\begin{align*}
  \sigma_{x,max}^{(1)} &= \frac{M_{b,max}^{(1)}}{I_z} \cdot e = \frac{M_{b,max}^{(1)}}{W_{z,min}} \\
  \sigma_{x,max}^{(1)} &= \sigma_{zul} = \frac{\sigma_F}{n} = \frac{355}{2} = 177,5 \, \text{MPa} \\
  \sigma_{zul} &= \frac{{M_{b,max}^{(1)}}}{W_{z,min}} \\
  &\boxed{
    W_{z,min} = 39,437 \, \text{cm}^3
  }
\end{align*}

\section{U-Profil auswählen}
Aus \href{https://www.mm.bme.hu/edu/bsc/de/tm2/DIN-1026-1-U-Profile.pdf}{der Tabelle} kann eine U-Profil mit dem kleinsten $W_y$ auswählen, wo $W_y > W_{min}$.
Die Profil mit Profilnummer \textbf{U240} habe ich gewählt, mit einer Widerstandsmoment $W_y = 39,6 \text{cm}^3$

\section{Berücksichtigung der Normalkraft und der Bohrung}
\subsection{Fläche und FTM des Profils}
Die Bohrung wird Material in dem gefährlichem Querschnitt abtrennen. Die Bohrung ist symmetrisch um die y-Achse, und geht durch eine Trapezförmige Querschnitt des Profils.

\begin{figure}[H]
  \centering
  \begin{tikzpicture}
    \draw[blue] (2.5, 2) -- (2.5, 4);
    \fill[pattern=horizontal lines, pattern color=blue] (2.5, 2) -- (3, 4) -- (2.5, 4);

    \draw[blue] (7.5, 2) -- (7, 0);
    \fill[pattern=horizontal lines, pattern color=blue] (7.5, 2) -- (7, 0) -- (7.5, 0);

    \draw[ultra thick] (0,0) -- (2,0) -- (3, 4) -- (0, 4) -- (0,0);
    \draw[ultra thick] (5,0) -- (7.5,0) -- (7.5, 4) -- (5, 4) -- (5,0);
  \end{tikzpicture}
  \caption{Vereinfachung des Profils}
  \label{fig:trapezinrechteck}
\end{figure}
Wenn die Biegung um die z-Achse passiert, kann das Querschnitt als im \figref{fig:trapezinrechteck} zu sehen ist vereinfacht werden. Dieser Vereinfachung ändert $I_z$ nicht.

Das Flächenträgheitmoment den gebohrten Querschnitte nach dem Formel:
\begin{align*}
  &I_{bohr} = 2 \cdot \frac{d^3 \cdot t}{12} = 138666,663 \, mm^4
\end{align*}

Die neue Fläche muss auch berechnet werden:
\begin{align*}
  &A = A_{U240} - A_{bohr} = 4230 - 2 \cdot 40 \cdot 13 = 3190 \, mm^2 \\
\end{align*}

Der neue Schwerpunkt kann von dem altem Schwerpunkt und der Schwerpunkt der Bohrung berechnet werden.
\begin{align*}
  s_z &= \frac{A_{U240} \cdot e_y - A_{bohr} \cdot \frac{b}{2}}{A} \\
  &= \frac{4230 \cdot 22{,}3 - 1040 \cdot 42{,}5}{3190} \\
  &= 15{,}71 \, \text{mm}
\end{align*}

Die Flächenträgheitmomente kann in diesem Schwerpunkt verschoben werden mit hilfe der Steiner-Satz.
\begin{align*}
  I_{S,bohr} &= I_{bohr} + A_{bohr} \cdot (\frac{b}{2} - s_z)^2 \\
  &= 885078{,}93 \, \text{mm}^4\\
  \\
  I_{S,U240} &= I_{U240} + A_{U240} \cdot (e_y - s_z)^2 \\
  &= 2480000 + 4230 \cdot (22{,}3 - 15{,}71)^2 \\
  &= 2663700{,}86 \, mm^4 \\
\end{align*}

Das Flächenträgheitmoment des neues Querschnitts kann jetzt berechnet werden.
\begin{align*}
  I_{z} &= I_{S,U240} - I_{S,bohr} \\
  I_{z} &= 2663700{,}86 - 885078{,}93 \\
  I_{z} &= 1778621{,}93 \, \text{mm}^4
\end{align*}

\subsection{Spannungen in dem Querschnitt}
Die Normalspannung wegen die Normalkraft:
\begin{align*}
  \sigma_{x,N} = \frac{N}{A} = \frac{10464 \, N}{3190 \, mm^2} = 3,28 \, \text{MPa}
\end{align*}

Die Normalspannung wegen der Biegemoment:
\begin{align*}
  \sigma_{x,b}(z) &= \frac{M_{b}^{(1)}(a)}{I_{z}} \cdot z \\
  &= \frac{7 \cdot 10^6 \, \text{Nmm}}{1778621{,}93 \, \text{mm}^4} \cdot z \\
  &= 3{,}936 \cdot z \, \text{[MPa]}
\end{align*}

Nach dem Superpositionsprinzip können die Zwei Spannungswerte addiert, so bekommt man die Spannunsverteilung in diesem Querschnitt:
\begin{align*}
  \sigma_{x}(z) &= \sigma_{x,N} + \sigma_{x,b}(z) \\
  &= 3{,}936z + 3,28 \, \text{[MPa]}
\end{align*}

\subsection{Dimensionierung}
Mit dem neuen Schwerpunkt ($s_z = 15,71$mm) ist die neue Randfase $e = -69,29$mm

Mit Substition bekommt die maximale Spannung in dem Querschnitt.
\begin{align*}
  \sigma_{x}(-69,29) &= -269,45 \, \text{MPa}
\end{align*}

Dieser Wert is deutlich größer, als die zulassige Spannung $\sigma_{zul} = 177,5$ MPa. Man muss also ein neues Profil ausgewählt werden.

\subsection{Neue Profil}
Da die U240 Profil nicht genug groß war, wird eine größere Profil ausgewählt, sei es die Profil \textbf{U280}.
De selben Berechnungen muss man machen als vorher, nur die Zahlenwerte des Profils sind anders.
\begin{align*}
  &A = A_{U280} - A_{bohr} = 5330 - 2 \cdot 40 \cdot 15 = 4130 mm^2 \\
  &I_{bohr} = 2 \cdot \frac{d^3 \cdot t}{12} = 160000 mm^4 \\
  &s_z = \frac{A_{U280} \cdot e_y - A_{bohr} \cdot \frac{b}{2}}{A} = \frac{5330 \cdot 25,3 - 1200 \cdot 47,5}{4130} = 18,85mm \\
  &I_{S,bohr} = I_{bohr} + A_{bohr} \cdot (\frac{b}{2} - s_z)^2 = 1144987 mm^4 \\
  &I_{S,U280} = I_{U280} + A_{U280} \cdot (e_y - s_z)^2 = 4211741,325 mm^4 \\
  &I_z = I_{S,U280} - I_{S,bohr} = 3066754,325 mm^4 \\
  &\sigma_{x,N} = \frac{N}{A} = \frac{10464 \text{N}}{4130 \text{mm}^2} = 2,5337 \, \text{MPa} \\
  &\sigma_{x,b}(z) = \frac{M_{b}^{(1)}(a)}{I_{z}} \cdot z = \frac{7 \cdot 10^6 \, \text{Nmm}}{3066754,325 \, \text{mm}^4} \cdot z = 2,2825 \cdot z \, \text{[MPa]} \\
  &\sigma_x(z) = 2,2825z + 2,5337 \, \text{[MPa]} \\
  &e = s_z - b = -76,15 mm \\
  &\sigma_x(-76,15) = -171,279 \text{MPa}
\end{align*}

Dieser Wert ist schon kleiner als die zulassige Spannung, also das Profil ist geeignet.

Zum Überprüfung wurde die \textbf{U260} Profil auch berechnet:
\begin{align*}
  &A = A_{U260} - A_{bohr} = 4830 - 2 \cdot 40 \cdot 14 = 3710 mm^2 \\
  &I_{bohr} = 2 \cdot \frac{d^3 \cdot t}{12} = 149333,33 mm^4 \\
  &s_z = \frac{A_{U260} \cdot e_y - A_{bohr} \cdot \frac{b}{2}}{A} = \frac{4830 \cdot 23,6 - 1120 \cdot 45}{3710} = 17,14mm \\
  &I_{S,bohr} = I_{bohr} + A_{bohr} \cdot (\frac{b}{2} - s_z)^2 = 1018654,482 mm^4 \\
  &I_{S,U260} = I_{U260} + A_{U260} \cdot (e_y - s_z)^2 = 3371563,628 mm^4 \\
  &I_z = I_{S,U260} - I_{S,bohr} = 2352909,146 mm^4 \\
  &\sigma_{x,N} = \frac{N}{A} = \frac{10464 \text{N}}{3710 \text{mm}^2} = 2,82 \, \text{MPa} \\
  &\sigma_{x,b}(z) = \frac{M_{b}^{(1)}(a)}{I_{z}} \cdot z = \frac{7 \cdot 10^6 \, \text{Nmm}}{2352909,146 \, \text{mm}^4} \cdot z = 2,975 \cdot z \, \text{[MPa]} \\
  &\sigma_x(z) = 2,975z + 2,82 \, \text{[MPa]} \\
  &e = s_z - b = -72,86 mm \\
  &\sigma_x(-72,86) = -213,9385 \text{MPa}
\end{align*}

Das U260 Profil hat eine größere Maximalspannung als die zulassige Spannung, also die U280 Profil ist die kleinste geeignete Profil.

\section{Schubspannungsverteilung}
\begin{figure}[H]
  \centering
  \begin{tikzpicture}
    \def\mmscale{0.05}
    \def\s{9.5 * \mmscale}
    \def\h{240 * \mmscale}
    \def\b{85 * \mmscale}
    \def\t{13 * \mmscale}

    \dimarrow{0,0}{0, -\b}{(-0.2,0)}{(-0.2,0)}{$b$}{left};
    \dimarrow{0,0}{\h, 0}{(0,0.2)}{(0,0.2)}{$h$}{above};
    \dimarrow{\h,0}{\h, -\s}{(0.2,0)}{(0.2,0)}{$s$}{right};
    \dimarrow{0,-\b}{\t, -\b}{(0,0.2)}{(0,0.2)}{$t$}{above};

    \draw[thick] (0,0) -- (\h, 0) -- (\h, -\s) -- (0, -\s) -- (0,0);
    \draw[thick] (0, -\s) -- (0, -\b) -- (\t, -\b) -- (\t, -\s);
    \draw[thick] (\h-\t, -\s) -- (\h-\t, -\b) -- (\h, -\b) -- (\h, -\s);

    \draw (\h / 2, -\s / 2) node[] {$A_1$};
    \draw (\t / 2, -\b / 2 - \s / 2) node[] {$A_2$};
    \draw (\h - \t / 2, -\b / 2 - \s / 2) node[] {$A_2$};
  \end{tikzpicture}
  \caption{Vereinfachter U240 Profil}
\end{figure}
Die Fläche und der Schwerpunkt muss berechnet werden.

\end{document}
